\documentclass[10pt,a4paper]{amsart}
\usepackage[margin=19mm]{geometry}
\usepackage{amsmath,amssymb,mathtools}
\usepackage[scr=rsfs,cal=euler]{mathalfa}
\usepackage{xfrac}
\usepackage[unicode,hidelinks,bookmarks=false]{hyperref}
\newcommand{\ie}{i.e.\ }
\newcommand{\cf}{cf.\ }
\newcommand{\resp}{respectively\ }
\newcommand{\Subsection}{\subsection}
\providecommand{\nobreakdash}{\nobreak\hbox{-}}
\setlength{\emergencystretch}{2em}
\multlinegap0pt
\usepackage{stmaryrd}
\usepackage{amssymb}
\usepackage{mathtools}
\usepackage{bm}
\newtheorem{thm}{Theorem}
\def\Z{\mathbb{Z}}
\def\ev{\text{ev.}}
\def\g{\mathfrak{g}}
\def\ua{\alpha}
\def\ub{\beta}
\title{Novel possible symmetries of $S$-matrix generated by $\Z_2^n$-graded Lie superalgebras}
\author{Ren Ito, Akio Nago}
\date{Source: arXiv:2501.07311v2 (CC BY 4.0)}
\begin{document}
\maketitle

\noindent\emph{Source and licence.} Novel possible symmetries of $S$-matrix generated by $\Z_2^n$-graded Lie superalgebras. Version arXiv:2501.07311v2 (CC BY 4.0), distributed by arXiv under the Creative Commons Attribution 4.0 International licence, http://creativecommons.org/licenses/by/4.0/, as recorded in the arXiv licence metadata for this exact version and shown on the abstract page https://arxiv.org/abs/2501.07311v2. Full licence notice: \texttt{licenses/arxiv-2501.07311-CC-BY-4.0.txt}.\par\smallskip

\noindent\textit{Editorial scope.} Counted unit: the article's thm thm (source0.tex lines 651--657). The article's complete original proof of it is delivered with it (source0.tex lines 658--670, one \texttt{proof} environment of 13 lines). No mathematics is written, repaired or re-derived: the statement and the proof are the article's own lines, in the article's own notation, and no retained line is substituted or re-flowed.  Retained uncounted as the source-local context the proof needs: lines 186--188; lines 640--643.  Nothing was omitted from any retained span, so the package carries no omission record.  No retained line contains a citation, so the wrapper carries no bibliography and no bibliography entry.  No retained line contains a figure environment, an image-inclusion command or drawing code, so no image is delivered and the package declares no asset dependency.  Editorial formatting only, in addition: an AMS \texttt{amsart} preamble that restates the article's own declarations and macros used by the retained lines, namely the environments \texttt{thm} on separate counters, together with the standard packages those lines need.  The retained environments print in this document as Theorem 1 in order of appearance, whereas the article prints the same items with the numbering of its own full document.  The complete native source of the article as distributed in the arXiv e-print for this exact version is bundled unchanged as \texttt{sources/171541/source0.tex}.
\begin{equation}
	\llbracket A, B \rrbracket = AB - (-1)^{\vec{a}\cdot\vec{b}} BA, \label{Z2n-LB}
\end{equation}

\begin{align}
 [A,B\}:=AB-(-1)^{\sum_{i=1}^n |\vec{\ua}_i\times \vec{\ub}_i|}BA,\quad A\in \g_{\vec{\ua}},\  B\in \g_{\vec{\ub}}
,\label{Z2n-LieB}
\end{align}

\begin{thm}
 The $(\Z_2^{2n+1})_{\ev}$-graded LSA 
\begin{align}
\g_{\ev}:=\bigoplus_{\vec{b}\in(\Z_2^{2n+1})_{\ev}}\g_{\vec{b}}
\end{align}
 is isomorphic to the $\Z_2^{2n}$-graded Lie algebra $\g_{LA}$.
\end{thm}

\begin{proof}
Let $f$ be an injective map $f: \vec{\ua}\in \Z_2^{2n}\to\vec{a}\in(\Z_2^{2n+1})_{\ev}$, defined by
\begin{align*}
\vec{a}=\left( \vec{\ua}_1, \vec{\ua}_2+|\vec{\ua}_1|(1,1),\cdots \vec{\ua}_i+\sum_{k=1}^{i-1}|\vec{\ua}_k|(1,1),\cdots, \vec{\ua}_{n}+\sum_{k=1}^{n-1}|\vec{\ua}_k|(1,1), \sum_{k=1}^{n}|\vec{\alpha}_k| \right)
\end{align*}
where  $|\vec{a}_k|^2:=a_{2k-1}^2+a_{2k}^2$, and it is equal to $|\vec{a}_k|=a_{2k-1}+a_{2k} \mod 2$.

Since $\Z_2^{2n}$-graded Lie algebra $\g_{LA}$ and $(\Z_2^{2n+1})_{\ev}$-graded LSA $\g_{\ev}$ have the same number of the graded subspaces, the only difference of them is the Lie bracket, which comes from $\sum_{i=1}^n |\vec{\ua}_i\times \vec{\ub}_i|$ in \eqref{Z2n-LieB} and $\vec{a}\cdot \vec{b}$ in \eqref{Z2n-LB}. However, taking into account that computations are performed in modulo 2, we have straightforwardly 
\begin{align}
 \sum_{i=1}^n |\vec{\ua}_i\times \vec{\ub}_i|=\vec{a}\cdot \vec{b}.
\end{align}
Therefore the map $f$ is an isomorphism between $\g_{LA}$ and $\g_{\ev}$.
\end{proof}

\end{document}
